


\chapter{Related Work}
\label{ch:related_work}
\section{Operational Technology Networks and Security Challenges}
\label{sec:ot_security_challenges}
Operational Technology (OT) refers to the hardware and software systems utilized for the direct monitoring, control, and actuation of physical devices, industrial processes, and events~\cite{kapoor_cyber_2025} They include components such as sensors, actuators, Programmable Logic Controllers (PLCs), Remote Terminal Units (RTUs), Supervisory Control and Data Acquisition (SCADA) systems, Human-Machine Interfaces (HMIs), and engineering workstations.  These systems form the basis of many Cyber-Physical Systems (CPS) and are used in critical infrastructure such as power grids, water treatment facilities, transportation networks and manufacturing plants~\cite{wang_guardians_2026}.

Historically, OT networks operated as isolated, air-gapped environments that relied almost entirely on physical separation for security~\cite{jeffrey_hybrid_2024}. This separation has changed with the adoption of Industry 4.0, an industrial paradigm that describes the digital transformation of manufacturing and other industrial sectors through connected, automated and data-driven systems. The Industrial Internet of Things (IIoT) is part of this shift and involves the connection of  industrial sensors, devices and machines to networks for data exchange, remote monitoring and automation. As a result, OT environments are increasingly integrated with traditional IT networks. While this IT-OT convergence improves data collection, automation and operational efficiency, it also expands the cyber-attack surface. Industrial assets that were once assumed to be isolated may become reachable from corporate networks, external services, or indirectly through compromised IT systems~\cite{kapoor_cyber_2025, jeffrey_hybrid_2024}.

Securing these environments is difficult because IT and OT networks do not always have the same operational priorities.  In traditional IT security, these priorities are often described using the CIA triad: Confidentiality, Integrity and Availability. Confidentiality protects information from unauthorized access, Integrity protects data and systems from unauthorized modification, and Availability ensures that systems remain accessible when needed. In traditional IT environmens, confidentiality often has the highest priority, followed by integraity and lastly availability. In OT environments, however, this ordering is reversed towards an AIC perspective, with availability being the main concern. This is because OT systems are directly connected to physical processes and are often subject to real-time and safety constraints. Downtime, delayed communication, or incorrect control actions can lead not only to data loss or service disruption, but also to production loss, equipment damage, environmental impact or risks to human safety~\cite{abshari_survey_2025, kapoor_cyber_2025}.



These risks have already been demonstrated by several well-known cyber-physical attacks. Stuxnet targeted Siemens PLCs used in an Iranian nuclear facility and is often considered one of the first examples of malware causing direct physical damage in an industrial process. It modified the control logic of uranium-enrichment centrifuges so that they operated outside safe conditions, while simultaneously hiding this abnormal behavior from monitoring systems by reporting normal-looking values~\cite{kapoor_cyber_2025, cabana_detecting_2019}. BlackEnergy and Industroyer later showed how attacks against power-grid environments can move beyond compromise of IT systems and directly affect electricity distribution. In the Ukrainian power-grid incidents, attackers interacted with industrial control systems and protocols to disrupt substations and cause real power outages~\cite{kapoor_cyber_2025, cabana_detecting_2019}. Triton/Trisis targeted Schneider Electric Triconex Safety Instrumented Systems in a petrochemical plant. This was especially concerning because safety instrumented systems are designed to keep industrial processes within safe operating limits, meaning that the attack targeted not only production systems but also the mechanisms intended to prevent catastrophic failures~\cite{kapoor_cyber_2025, wang_guardians_2026}. More recent ransomware incidents, such as EKANS and Colonial Pipeline, further show that cyber incidents do not need to directly manipulate PLC logic to affect industrial operations. EKANS included functionality to terminate industrial control processes, while the Colonial Pipeline incident led to a multi-day shutdown of fuel pipeline operations~\cite{kapoor_cyber_2025,  ding_data-driven_2022}. This illustrates why OT security has to consider cyber-physical consequences, not only traditional information security risks.

Furthermore, the deployment of security solutions in OT is often limited by legacy systems and resource constraints. Many OT environments still rely on older devices and industrial protocols, such as Modbus, DNP3 and EtherNet/IP, which were designed mainly for reliable communication and not necessarily with built-in security mechanisms such as encryption or strong authentication~\cite{wang_guardians_2026, kapoor_cyber_2025}. Updating or replacing such systems can be operationally risky, expensive, or delayed because of uptime requirements~\cite{abshari_survey_2025, kapoor_cyber_2025}. In addition, endpoint devices at the lower levels of an ICS architecture, such as sensors, actuators and PLCs, often have limited memory and computational resources~\cite{abshari_survey_2025, jeffrey_hybrid_2024}. This makes it difficult to deploy heavy host-based security mechanisms directly on these devices.


For these reasons, it is important to explore security approaches that are adapted to OT environments. Any monitoring or detection method has to take into account the operational constraints of industrial systems, where availability, safety and continuity are central requirements. At the same time, the method should be deployable without major changes to legacy devices or interruptions of the physical process. This is a difficult balance. OT environments need better visibility into suspicious behavior, but the security mechanism itself should not become a source of operational risk. The next section therefore discusses anomaly detection as a possible approach for detecting abnormal behavior in OT networks.


\section{Anomaly Detection in OT Environments}
Given the constraints discussed in Section~\ref{sec:ot_security_challenges}, OT security mechanisms should provide visibility into suspicious behavior without disrupting the underlying physical process. Intrusion Detection Systems (IDS) are one way to provide this visibility and are commonly divided into signature-based and anomaly-based approaches. Signature-based IDS compare observed traffic or system events against predefined rules or patterns, known as signatures, that describe previously identified malicious activity. This makes them effective for attacks that have already been characterized, but less suitable for novel or targeted attacks whose signatures are not yet availabl~\cite{meira_performance_2020, jeffrey_hybrid_2024} This limitation is especially relevant in OT environments, where attacks may be highly system-specific and where generating signatures for complex process or network behavior can be difficult~\cite{wang_guardians_2026}.

Anomaly-based detection takes a different approach: instead of searching for known attack patterns, normal behavior is modeled based on observed data, and deviations from this model are flagged as suspicious~\cite{abshari_survey_2025}. Anomaly detection has been studied extensively in IT security, but methods developed for general IT traffic cannot simply be assumed to transfer directly to OT environments. IT traffic is often shaped by human activity, web applications, cloud services and changing user behavior, making it diverse and highly variable. OT traffic, in contrast, is more closely tied to the physical process and is often driven by fixed interactions between devices, periodic control tasks, polling behavior and stable asset roles. This makes anomaly detection attractive for OT networks, where unknown attacks are an important concern and where communication is often more regular than in general IT networks. This periodic nature of OT networks, means deviations from expected communication can provide more useful indicators of abnormal behavior.~\cite{meira_performance_2020, catillo_cps-guard_2023}. 

However, this regularity does not mean that a single generic baseline can be applied to all OT networks. Jeffrey et al.~\cite{jeffrey_hybrid_2024} describe IT networks as being closer to a ``monoculture'', with many similar systems and more standardized communication protocols, while OT networks are much more heterogeneous and differ strongly between organizations. A pattern that is normal in one industrial environment may therefore be unusual in another. For this reason, OT anomaly detection cannot simply reuse IT-based methods or generic traffic models, but requires environment-specific baselines that capture the temporal and asset-level behavior of the monitored operational context~\cite{jeffrey_hybrid_2024}.



A major challenge in developing machine learning-based anomaly detection for OT environments is the limited availability of labeled attack data. Cyberattacks on critical infrastructure are rare, difficult to observe, and unsafe to reproduce in operational environments. As a result, real-world OT datasets are often dominated by benign traffic, while labeled malicious examples are scarce or unavailable~\cite{jeffrey_hybrid_2024, wang_guardians_2026}. This makes fully supervised learning difficult, since such models require representative examples of both normal and attack behavior. Even when attack labels are available in public datasets, they may not cover the full range of realistic attacks or environment-specific behavior. For this reason, one-class and semi-supervised approaches are frequently used in OT anomaly detection~\cite{catillo_cps-guard_2023, jeffrey_hybrid_2024, meira_performance_2020}. In this setting, the model is trained mainly or exclusively on benign data and learns a representation of normal operation. During detection, samples that deviate sufficiently from this learned normal behavior receive a high anomaly score.

\section{Limitations of Existing OT Anomaly Detection}

Although one-class and semi-supervised anomaly detection methods are a practical direction for OT networks, existing research still has several limitations that make it difficult to transfer results to real operational environments. A first limitation is the strong reliance on a small number of public benchmark datasets, such as SWaT, WADI and BATADAL~\cite{atluri_perspectives_2022, abshari_survey_2025}. These datasets are valuable because they provide labeled attack scenarios and allow different methods to be tested on a common baseline. However, they are usually collected in controlled test environments and therefore do not fully represent the conditions found in real OT networks. Operational environments can contain noisy traffic, benign misconfigurations, maintenance activity, incomplete labels or customer-specific behavior that are difficult to capture in benchmark datasets~\cite{moriano_adaptive_2025, jeffrey_hybrid_2024}. As a result, models that perform well on clean and controlled datasets may not generalize directly to real, noisy environments.

Another limitation is that results across OT anomaly detection studies are difficult to compare directly. Existing work evaluates a wide range of methods, from distance-based and clustering approaches to tree-based outlier detectors and deep reconstruction-based models, but these are often tested with different datasets, preprocessing steps, feature sets, thresholding strategies and model configurations~\cite{wang_guardians_2026, atluri_perspectives_2022}. As a result, it is difficult to determine whether a reported improvement is caused by the model architecture itself, by the data pipeline, or by the evaluation setup. If models cannot be compared under consistent conditions, it is also difficult to decide which approach is actually suitable for a specific OT monitoring problem. Recent evaluations further show that preprocessing choices and detection hyperparameters can strongly influence the final result, while the specific model architecture is not always the dominant factor~\cite{atluri_perspectives_2022}. Hence comparing candidate models within the same data pipeline and evaluation setup leads to more useful conclusions rather than treating reported performance numbers from different studies as directly comparable.

A final limitation of much OT anomaly detection research is that it focuses mainly on physical process variables, such as sensor readings, actuator states, tank levels or other process measurements. This is natural for datasets such as SWaT and WADI, where attacks are often represented through changes in the physical process. However, it also means that many proposed methods are evaluated on attacks that have already affected the process itself~\cite{atluri_perspectives_2022, behdadnia_early-stage_2025}. Less attention is given to earlier network-level activity, such as host discovery, port scanning or other reconnaissance behaviour, even though these actions may occur before direct manipulation of the physical infrastructure.



\section{Early-Stage Detection and Reconnaissance}
\label{sec:early_stage_detection}

The focus on physical process variables can be understood through the Cyber Kill Chain (CKC), which as depicted in Figure~\ref{fig:cyber_kill_chain}, describes an intrusion as a sequence of stages: reconnaissance, weaponization, delivery, exploitation, installation, command and control, and action on objectives~\cite{behdadnia_early-stage_2025}. In the early stages, the attacker gathers information about the target environment, prepares or delivers the attack, and attempts to gain a foothold. In the later stages, the attacker establishes persistence, communicates with compromised systems, and ultimately executes their objective. In an OT context, this final objective may involve manipulating the physical process, for example through false data or command injection~\cite{behdadnia_early-stage_2025}. Process-level anomaly detection mainly observes these later stages, where the attacker has already reached the point of affecting the physical infrastructure. While such detection is still important, it leaves less time for operators to react before operational consequences occur.

For this reason, early-stage detection has become an important direction for OT security. The initial stages of an intrusion can last much longer than the final execution phase, providing defenders with more time to detect suspicious behavior and respond before the attack reaches its objective~\cite{behdadnia_early-stage_2025}

Reconnaissance is especially relevant in this context because it is often one of the first steps performed by an attacker. Before exploiting a target, the attacker typically tries to understand the network by identifying active hosts, open ports, exposed services and possible entry points. This can involve host discovery, port scanning and service enumeration~\cite{keshavamurthy_early_2023, cabana_detecting_2019}. These activities cannot be observed through physical processes but leave observable traces in network communication, such as connections to many destinations, probes to many ports, short-lived TCP flows or failed connection attempts. Detecting such patterns can provide an early warning before the attacker moves further along the kill chain~\cite{zaidy_comprehensive_2025, ring_detection_2018}.



% \begin{figure}[htbp]
%     \centering
%     \makebox[\textwidth][c]{%
%     \begin{adjustbox}{width=1.3\textwidth}
%     \begin{tikzpicture}[
%         tactic/.style={
%             rectangle,
%             rounded corners=2pt,
%             draw=black!60,
%             minimum width=2.15cm,
%             minimum height=0.75cm,
%             align=center,
%             font=\footnotesize,
%             text width=1.75cm
%         },
%         normal/.style={tactic, fill=blue!10},
%         focus/.style={tactic, fill=green!18},
%         impact/.style={tactic, fill=red!18}
%     ]

%     \node[normal] (initial) {Initial\\Access};
%     \node[normal, right=0.20cm of initial] (execution) {Execution};
%     \node[normal, right=0.20cm of execution] (persistence) {Persistence};
%     \node[normal, right=0.20cm of persistence] (privilege) {Privilege\\Escalation};
%     \node[normal, right=0.20cm of privilege] (evasion) {Evasion};
%     \node[focus, right=0.20cm of evasion] (discovery) {Discovery};
%     \node[normal, right=0.20cm of discovery] (lateral) {Lateral\\Movement};
%     \node[normal, right=0.20cm of lateral] (collection) {Collection};
%     \node[normal, right=0.20cm of collection] (c2) {Command\\and Control};
%     \node[impact, right=0.20cm of c2] (inhibit) {Inhibit\\Response\\Function};
%     \node[impact, right=0.20cm of inhibit] (impair) {Impair\\Process\\Control};
%     \node[impact, right=0.20cm of impair] (impactnode) {Impact};

%     \node[above=0.18cm of discovery, font=\footnotesize\bfseries, text=green!45!black]
%         {Focus of this thesis};

%     \node[below=0.18cm of impair, font=\footnotesize\bfseries, text=red!60!black]
%         {Physical/process impact};

%     \end{tikzpicture}
%     \end{adjustbox}%
%     }
%     \caption{MITRE ATT\&CK for ICS tactics.}
%     \label{fig:mitre_attack_ics_tactics}
% \end{figure}

\begin{figure}[htbp]
    \centering
    \makebox[\textwidth][c]{%
    \begin{adjustbox}{width=1.3\textwidth}
    \begin{tikzpicture}[
        tactic/.style={
            rectangle,
            rounded corners=2pt,
            draw=black!60,
            minimum width=2.15cm,
            minimum height=0.75cm,
            align=center,
            font=\footnotesize,
            text width=1.75cm
        },
        normal/.style={tactic, fill=blue!10},
        focus/.style={tactic, fill=green!18},
        impact/.style={tactic, fill=red!18},
        techlabel/.style={
            anchor=west,
            font=\scriptsize,
            align=left,
            text=black!85
        }
    ]

    % Main tactic row
    \node[normal] (initial) {Initial\\Access};
    \node[normal, right=0.20cm of initial] (execution) {Execution};
    \node[normal, right=0.20cm of execution] (persistence) {Persistence};
    \node[normal, right=0.20cm of persistence] (privilege) {Privilege\\Escalation};
    \node[normal, right=0.20cm of privilege] (evasion) {Evasion};
    \node[focus, right=0.20cm of evasion] (discovery) {Discovery};
    \node[normal, right=0.20cm of discovery] (lateral) {Lateral\\Movement};
    \node[normal, right=0.20cm of lateral] (collection) {Collection};
    \node[normal, right=0.20cm of collection] (c2) {Command\\and Control};
    \node[impact, right=0.20cm of c2] (inhibit) {Inhibit\\Response\\Function};
    \node[impact, right=0.20cm of inhibit] (impair) {Impair\\Process\\Control};
    \node[impact, right=0.20cm of impair] (impactnode) {Impact};

    % Labels
    \node[above=0.18cm of discovery, font=\footnotesize\bfseries, text=green!45!black]
        {Focus of this thesis};

    \node[below=0.18cm of impair, font=\footnotesize\bfseries, text=red!60!black]
        {Physical/process impact};

    % ------------------------------------------------------------------
    % Discovery branch: moved toward the LEFT side of the Discovery box
    % ------------------------------------------------------------------

    % Start near left side of Discovery box, not in the center
    \coordinate (discstart) at ($(discovery.south west)+(0.22cm,-0.05cm)$);

    % First level: Remote System Discovery
    \coordinate (remotelevel) at ($(discstart)+(0,-0.42cm)$);
    \draw[thick, green!45!black] (discstart) -- (remotelevel);
    \draw[thick, green!45!black] (remotelevel) -- ++(0.42cm,0) coordinate (remoteend);
    \node[techlabel] at ($(remoteend)+(0.10cm,0)$)
        {Remote System Discovery (T0846)};

    % Second level: Port Scan as sub-technique of Remote System Discovery
    \coordinate (subvertend) at ($(remoteend)+(0,-0.42cm)$);
    \draw[thick, green!45!black] (remoteend) -- (subvertend);
    \draw[thick, green!45!black] (subvertend) -- ++(0.35cm,0) coordinate (portend);
    \node[techlabel] at ($(portend)+(0.10cm,0)$)
        {Port Scan (T0846.001)};

    \end{tikzpicture}
    \end{adjustbox}%
    }
    \caption{MITRE ATT\&CK for ICS tactics and the Discovery-related behaviour studied in this thesis.}
    \label{fig:mitre_attack_ics_tactics}
\end{figure}

\section{Unsupervised Models for Network-Based Anomaly Detection}
\label{sec:literature_models}
To detect the early-stage reconnaissance activities discussed in Section~\ref{sec:early_stage_detection}, multiple unsupervised and semi-unsupervised models were explored that can identify deviations in network traffic without relying on known attack signatures. 

Distance-based algorithms, and especially k-Nearest Neighbours (k-NN), are frequently used as lightweight anomaly detection models~\cite{meira_performance_2020, jeffrey_hybrid_2024, cai_adam_2023}. In a one-class setting, k-NN stores benign reference samples and scores a new observation according to its distance from the closest benign samples.  Meira et al.~\cite{meira_performance_2020} evaluated several unsupervised techniques for cyber-attack anomaly detection on the public NSL-KDD and ISCX network intrusion datasets, including One-Class Nearest Neighbour, One-Class K-Means, Isolation Forest, One-Class SVM and an autoencoder. Their results show that nearest-neighbour methods can perform competitively on network intrusion data, especially on the ISCX dataset~\cite{meira_performance_2020}. Jeffrey et al. also motivate the use of one-class models in CPS environments, arguing that malicious data is often scarce and that modelling benign behaviour directly is more suitable than relying on balanced two-class training data~\cite{jeffrey_hybrid_2024}. In their experiments, one-class k-NN achieved the best accuracy among the evaluated one-class and two-class models, which supports its use as a simple but strong baseline~\cite{jeffrey_hybrid_2024}.

A related line of work combines distance-based detection with statistical descriptions of network behavior. Entropy-based features are particularly relevant for network anomaly detection because they summarize how concentrated or dispersed communication is across hosts, ports or protocols. The ADAM scheme proposed by Cai et al.~\cite{cai_adam_2023} uses entropy vectors as traffic profiles and combines them with unsupervised anomaly detection for Distributed Denial of Service (DDoS) mitigation in software-defined cyber-physical systems. In particular, ADAM selects k-NN as the anomaly detection module because it provides a practical balance between detection effectiveness and the runtime overhead of classifying new entropy vectors. Although ADAM targets DDoS attacks rather than reconnaissance, it is relevant for this thesis because reconnaissance can also change the distribution of contacted destinations and ports. This motivates the inclusion of entropy-based feature groups when evaluating distance-based models.

Clustering approaches have also shown promising results in industrial anomaly detection. For one-class anomaly detection, K-Means can represent benign behavior as a set of clusters, after which new samples are scored by their distance to the nearest cluster centroid. Wang et al.~\cite{wang_guardians_2026} compare several supervised and unsupervised models on the SWaT testbed and report that K-Means performs strongly among the unsupervised methods, achieving the highest recall and F1-score among the unsupervised techniques and one of the lowest undetected rates. In this context, recall measures how many attacks are detected, while F1-score balances detection recall with false alarms. This suggests that relatively simple clustering models can capture useful structure in deterministic OT environments. However, this result is still based on a controlled benchmark testbed. It therefore remains useful to evaluate whether the same type of simple clustering model can remain effective on noisier operational network metadata.

Alongside distance-based and clustering methods, reconstruction-based deep learning models are widely used in CPS and OT anomaly detection. Autoencoders are trained to reconstruct their input after compressing it through a lower-dimensional representation. When trained on benign data, the underlying assumption is that normal samples will be reconstructed well, while anomalous samples will result in a higher reconstruction error. Ortega-Fernandez et al.~\cite{ortega-fernandez_network_2024} propose a deep-autoencoder-based network intrusion detection system for DDoS attacks in ICS using network flow data. Their work is especially relevant because it focuses on flow-based features rather than device logs or process variables, requires limited preprocessing and evaluates the feasibility of the approach on data of a real industrial plant. However, for safety reasons, attacks could not be executed in the production environment, so the real-world evaluation mainly concerns false alarms and deployment feasibility, while attack detection is evaluated on a separate ICS DDoS dataset. 

A practical challenge for autoencoder-based methods is that benign training data is not always perfectly clean. Operational data may contain rare benign behavior, measurement noise, imperfect labels or hidden outliers. Catillo et al.~\cite{catillo_cps-guard_2023} address this issue in CPS-GUARD, which combines a semi-supervised autoencoder with an outlier-aware threshold selection strategy based on Isolation Forest. With a focus on practical deployability, they argue that CPS intrusion detection should favor semi-supervised learning, non-intrusive data sources and simple model designs over unnecessarily complex architectures.

Finally, recent evaluations suggest that increasing model complexity does not automatically improve anomaly detection performance. Fung et al. perform a comprehensive evaluation of reconstruction-based anomaly detection in ICS and show that model architecture and size are not always the dominant factors, with smaller models often performing similarly to larger architectures~\cite{fung_perspectives_2022}. This supports the use of simple models as meaningful baselines rather than treating them only as weak alternatives to complex deep learning models.


\section{Evaluation Metrics}
The choice of evaluation metric is an important part of anomaly detection research. When malign data is sparse, overall accuracy can be misleading, as a model  that predicts almost everything as benign may still obtain a high accuracy while failing to detect the relevant attacks~\cite{wang_guardians_2026, fung_perspectives_2022}. For this reason, anomaly detection studies commonly report precision, recall and F1-score, which provide a more informative view of the trade-off between false positives and false negatives.

However, standard precision, recall and F1-score are usually computed point-wise, meaning that each timestamp or window is evaluated independently. This is not always aligned with the way attacks appear in practice. Many OT and network attacks unfold over a temporal interval rather than as isolated samples. A reconnaissance scan, for example, may generate a sequence of anomalous windows, while a low-and-slow scan may be spread over a much longer period. Point-wise metrics do not explicitly distinguish between detecting an attack immediately and detecting it near the end of the attack interval. They can also give disproportionate weight to long attacks, because every anomalous window contributes separately to the final score~\cite{fung_perspectives_2022}.


To address this, recent work has proposed the use of range-based metrics for temporal anomaly detection. Instead of evaluating each window independently, range-based metrics group consecutive anomalous windows into temporal ranges and compare predicted anomaly ranges with the true attack ranges~\cite{fung_perspectives_2022}. This better matches operational alert analysis, where the main concern is whether an attack episode is detected, how fragmented the alerts are, and whether detection occurs early enough to support response. This distinction matters because the metric used for tuning and evaluation can change which model appears best. Fung et al.~\cite{fung_perspectives_2022} show that evaluating ICS anomaly detectors with range-based metrics can lead to different conclusions than using point-F1 alone.


\section{Thesis Positioning}
The reviewed literature motivates three main choices in this thesis. First, the detection problem is studied using operational network metadata rather than clean benchmark data or physical process variables. The input consists of Linux conntrack metadata collected from AXS Guard appliances, which provides lightweight visibility into communication behaviour at the network edge without relying on packet payloads or host-based agents.

Second, the thesis focuses on early-stage reconnaissance instead of later-stage physical process manipulation. Host discovery, port scanning, service detection and low-and-slow scanning are treated as deviations from the normal communication behaviour of a source host. This places the work earlier in the Cyber Kill Chain, where detection can provide more time for investigation before the physical process is affected.

Third, the thesis evaluates existing unsupervised model families under a common experimental setup. Instead of proposing a new model architecture, it compares lightweight distance- and clustering-based methods with a reconstruction-based autoencoder approach using the same windowing, feature extraction, thresholding and evaluation pipeline. This makes it possible to study whether models that perform well in the literature remain useful on noisy operational OT metadata.


The following chapters therefore examine how conntrack flows are reconstructed, how reconnaissance attacks are simulated and injected into real benign traffic, and how window-level features and anomaly detection models are evaluated using both point-based and range-based metrics.


% \section{Positioning of This Thesis}
% concise bridge to your dataset/methodology chapters.